\section{Applications}
\label{sec:apps}

We wrote our phone recordings around names, numbers, drug and course terms, and the clips in Table~\ref{tab:examples} point to settings where an on-device, code-switching recogniser is useful and where the embedded word carries the content:
\begin{itemize}
  \item \textbf{Clinics and pharmacies.} Notes or patient instructions such as ``\zh{医生说} prednisolone \zh{要慢慢减量}'' or ``the nurse said \zh{布洛芬} only after meals''. Patient audio should stay on the device, and the drug name is exactly the word a wrong transcript must not lose.
  \item \textbf{Live captions for deaf and hard-of-hearing users} in lectures, tutorials and meetings, where course codes and technical terms (module, quiz, release) are embedded in Mandarin speech. Captions must keep up with speech: small FT+L in 8-bit runs at RTF 0.57 on a laptop CPU.
  \item \textbf{Banking, government and telco hotlines}, where callers mix in words like \zh{身份证} or \zh{贷款} over a phone line. Our phone recordings show that off-the-shelf Whisper writes these as English sound-alikes, while the fine-tuned models get them right.
  \item \textbf{Eldercare and family use}, e.g.\ voice reminders like ``my mother keeps the \zh{药} in the second drawer''. A home device should not stream household speech to the cloud, and it must still work without a fast network connection.
  \item \textbf{Personal dictation} of messages, notes and reminders (``the \zh{期末考} starts on the first of December''), where a missed embedded word loses the one fact the note was written to keep.
\end{itemize}

\paragraph{Conversational use.} The same recogniser can be the input to a spoken assistant: the user speaks, the on-device model transcribes each turn, a language model reads the transcript and writes a reply, and a speech synthesiser reads that reply aloud. In such a dialogue the embedded word is often the slot the assistant must act on (the drug, the date, the room), so the role-based error we measure matters more than overall MER. A dropped embedded word does not produce an obviously broken transcript. It produces a fluent sentence that the language model will answer confidently and wrongly (``bring your shuffling jump'' instead of \zh{身份证}). Our language-ID head can also help here: it tells the assistant which language each part of the turn was in, so the reply can follow the user's mix, and, because our best system still misses about one embedded word in five on NSC, the assistant can read back an embedded word before acting on it (``Prednisolone, is that right?'').

Turn-taking sets the speed budget. At RTF 0.57, a 5-second turn takes about 3\,s to transcribe on our laptop, which is too slow for a natural reply; base at RTF 0.21 (CTranslate2) takes about 1\,s, and gated re-ranking adds cost only on the quarter of turns with a switch. A conversational system would therefore pair a smaller or streaming recogniser with our loss and gating methods, and could pass the previous turns to Whisper as a text prompt so that terms already mentioned in the dialogue are more likely to be recognised again. We have not tested these designs; they are directions this work makes possible, not results.
